\subsection{Двумерная гауссиана}

\subsubsection{Алгоритм}

\textit{Двумерная гауссиана} --- это нелинейная функция следующего вида:
\[
z(x, y) = A \cdot \exp \left( - \left[ \frac{\tilde{x}^2}{2\sigma_x^2} + \frac{\tilde{y}^2}{2\sigma_y^2} \right] \right) + z_0
\]
где:
\[
\begin{pmatrix}
\tilde{x} \\
\tilde{y}
\end{pmatrix}
=
\begin{pmatrix}
\cos\theta & \sin\theta \\
-\sin\theta & \cos\theta
\end{pmatrix}
\cdot
\begin{pmatrix}
x - x_0 \\
y - y_0
\end{pmatrix}
\]

Параметры функции:
\begin{itemize}
  \item $A$ --- амплитуда \textit{(высота пика относительно уровня смещения)};
  \item $x_0, y_0$ --- координаты центра функции;
  \item $\sigma_x, \sigma_y$ --- среднеквадратичные отклонения вдоль повёрнутых осей;
  \item $\theta$ --- угол поворота гауссианы относительно осей $x$ и $y$;
  \item $z_0$ --- константное смещение поверхности.
\end{itemize}

Имеем задачу регрессии, которая заключается в минимизации среднеквадратичной ошибки:
\[
\begin{gathered}
  Loss = L(\{x_i,y_i,z_i\}^N_{i=1},A,z_0,\sigma_x,\sigma_y,x_0,y_0,\theta) = MSE = \\
  =\frac{1}{2\cdot N} \sum_{i=1}^{N} (z_i - z(x_i,y_i))^2 \to \min_{A,x_0,y_0,\sigma_x,\sigma_y,z_0,\theta}
\end{gathered}
\]

Можно заметить, что некоторые параметры модели находятся в показателе экспоненты. Из этого следует, что применить линейную регрессию не получится. Свести функцию к обобщённой линейной регрессии тоже не выйдет из-за параметра \( z_0 \): если бы его не было, то можно прологарифмировать обе части и получить обобщённую линейную регрессию.

Для нелинейной регрессии отлично подойдёт квазиньютоновский метод L-BFGS-B:

\subsubsection{Кривая обучения}

\dots

\subsubsection{Вектор невязки}

\dots

\subsubsection{Модельная поверхность}

\dots